\documentclass[11pt,a4paper]{article}
\usepackage[margin=2.5cm]{geometry}
\usepackage{amsmath,amsthm}
\usepackage{newtxtext}
\usepackage{newtxmath}
\usepackage{enumitem}
\usepackage{microtype}
\setlength{\parindent}{0pt}
\setlength{\parskip}{4pt}
\allowdisplaybreaks

\newcommand{\E}{\mathbb{E}}
\renewcommand{\P}{\mathbb{P}}
\newcommand{\R}{\mathbb{R}}
\newcommand{\N}{\mathbb{N}}
\newcommand{\Q}{\mathbb{Q}}
\newcommand{\F}{\mathcal{F}}
\newcommand{\G}{\mathcal{G}}
\newcommand{\Hc}{\mathcal{H}}
\newcommand{\1}{\mathbf{1}}
\newcommand{\Leb}{\mathrm{Leb}}
\newcommand{\sgn}{\operatorname{sgn}}
\newcommand{\prob}[1]{\par\medskip\noindent\textbf{#1}\quad}
\newcommand{\sub}[1]{\par\smallskip\noindent\textbf{#1}\ }
\newcommand{\note}[1]{\par\noindent\emph{Remark.} #1}

\begin{document}

\begin{center}
{\large\textbf{Stochastic Processes --- 1060429}}\\[2pt]
Nawar Sgier, 326586047
\end{center}

\noindent Throughout, a Brownian motion (BM) is standard ($B_0=0$) with continuous paths, and
$\F_t$ denotes the filtration in use. Apparent typos are pointed out where they occur.

%=====================================================================
\section*{Exercise 1}
\subsection*{Relations between stochastic processes}

\prob{1.}
Let $\Omega_0=\{X,Y\text{ right-continuous}\}\cap\bigcap_{q\in\Q_+}\{X_q=Y_q\}$. The first event
contains a measurable set of probability $1$, and each $\{X_q=Y_q\}$ has probability $1$
(modification), so $\Omega_0$ contains a measurable set of probability $1$ (countable
intersection). For $\omega\in\Omega_0$ and $t\ge0$ take $q_n\in\Q_+$, $q_n\downarrow t$; then
$X_t(\omega)=\lim X_{q_n}(\omega)=\lim Y_{q_n}(\omega)=Y_t(\omega)$. Hence $X_t=Y_t$ for all $t$ on
$\Omega_0$, i.e.\ $X,Y$ are indistinguishable.

\prob{2.}
Let $Q_0=\Q\cap[0,t_0)$ and
\[
 C_{t_0}=\bigcap_{k\ge1}\bigcup_{m\ge1}\ \bigcap_{s,u\in Q_0,\ |s-u|<1/m}\{|X_s-X_u|\le 1/k\}
 \in\sigma(X_t:t\le t_0).
\]
\textit{Claim.} If the path $X(\omega)$ is RCLL on $[0,t_0]$, then $\omega\in A_{t_0}\iff \omega\in C_{t_0}$.

If $X(\omega)$ is continuous on $[0,t_0)$ and has a left limit at $t_0$, it extends continuously to
$[0,t_0]$, so it is uniformly continuous on $[0,t_0)$ and $\omega\in C_{t_0}$. Conversely, let
$\omega\in C_{t_0}$, $k\ge1$, and $m$ as in the definition. For $s,u\in[0,t_0)$ with $|s-u|<1/m$
choose $q_j,r_j\in Q_0$ with $q_j\downarrow s$, $r_j\downarrow u$ and $|q_j-r_j|<1/m$ (possible as
$s,u<t_0$). By right-continuity $|X_s-X_u|=\lim|X_{q_j}-X_{r_j}|\le1/k$. So $X(\omega)$ is
uniformly continuous on $[0,t_0)$, hence $\omega\in A_{t_0}$.

\sub{(a)} If every path is RCLL, the claim gives $A_{t_0}=C_{t_0}\in\sigma(X_t:t\le t_0)$.

\sub{(b)} Let $\Omega=[0,1]$, $\F=\mathcal B([0,1])$, $\P=\Leb$, and let $C$ be the Cantor set. Since
$|\mathcal P(C)|=2^{\mathfrak c}>\mathfrak c=|\mathcal B([0,1])|$, there is $E\subset C$ with
$E\notin\F$. Define $X_t(\omega)=\1_{E}(\omega)\1_{\{t=\omega\}}$. Then $X_t=\1_{E\cap\{t\}}$ is the
indicator of the empty set or of a single point, so it is measurable. For $\omega\notin E$ the path is $\equiv0$. For
$\omega\in E$ the path has a spike at $t=\omega$ and is not right-continuous there. Thus
$\{X\text{ RCLL}\}=\Omega\setminus E\supseteq\Omega\setminus C$ and $\P(\Omega\setminus C)=1$, so
$X$ is RCLL with probability $1$. But $A_2=\Omega\setminus E\notin\F$, and in particular
$A_2\notin\sigma(X_t:t\le 2)\subseteq\F$.

\sub{(c)} There is $N\in\F$ with $\P(N)=0$ such that every path outside $N$ is RCLL. By the claim,
$A_{t_0}\triangle C_{t_0}\subseteq N$. Completeness of $\F_{t_0}$ (it contains every subset of a
$\P$-null set) gives $A_{t_0}\setminus C_{t_0},\,C_{t_0}\setminus A_{t_0}\in\F_{t_0}$. Hence
$A_{t_0}=(C_{t_0}\setminus(C_{t_0}\setminus A_{t_0}))\cup(A_{t_0}\setminus C_{t_0})\in\F_{t_0}$.

\subsection*{Brownian motion}

\prob{1.}
\note{The condition should read $\E[|X_t-X_s|^\alpha]\le M|t-s|^{1+\beta}$ for all $t,s\in[0,T]$. In (c) the
bound should be $2^{-\gamma n}$, not $2^{-m}$. In (e),(f) the index set is $[0,1]$.}

\sub{(a)} Put $\tilde X_u=X_{uT}$, $u\in[0,1]$. Then $\E|\tilde X_u-\tilde X_v|^\alpha\le
MT^{1+\beta}|u-v|^{1+\beta}$, which is the hypothesis on $[0,1]$ with constant $MT^{1+\beta}$. If
$\tilde Y$ is a modification of $\tilde X$ that is a.s.\ $\gamma$-H\"older, then
$Y_t=\tilde Y_{t/T}$ satisfies $\P(Y_t=X_t)=\P(\tilde Y_{t/T}=\tilde X_{t/T})=1$, and
$|Y_t-Y_s|\le K|t/T-s/T|^\gamma=KT^{-\gamma}|t-s|^\gamma$.

\sub{(b)} By Markov's inequality, $\P(|X_t-X_s|>\varepsilon)\le\varepsilon^{-\alpha}M|t-s|^{1+\beta}\to0$ as $t\to s$.

\sub{(c)} By a union bound and Markov's inequality,
\[
\P\Big(\max_{1\le k\le 2^n}|X_{k2^{-n}}-X_{(k-1)2^{-n}}|>2^{-\gamma n}\Big)
\le 2^n\cdot 2^{\gamma\alpha n}M2^{-n(1+\beta)}=M2^{-n(\beta-\alpha\gamma)},
\]
which is summable since $\gamma<\beta/\alpha$. By Borel--Cantelli, a.s.\ only finitely many of these
events occur. Let $N(\omega)$ be the last $n$ at which one occurs ($N=0$ if none occurs). Then
$\P(N<\infty)=1$ and $\max_k|X_{k2^{-n}}-X_{(k-1)2^{-n}}|\le2^{-\gamma n}$ for all $n>N(\omega)$.

\sub{(d)} Fix $\omega$ with $N(\omega)<\infty$.
\emph{Claim:} for $m>n\ge N(\omega)$ and $s<t$ in $D_m$ with $t-s<2^{-n}$,
\[|X_t-X_s|\le2\sum_{j=n+1}^m2^{-\gamma j}.\]
We use induction on $m$. For $m=n+1$ we have $t-s=2^{-(n+1)}$, and the bound holds by (c). For the step
$m\to m+1$, let $s,t\in D_{m+1}$. Let $s'=\min\{r\in D_m:r\ge s\}$ and $t'=\max\{r\in D_m:r\le t\}$. Then
$s'-s,\ t-t'\in\{0,2^{-(m+1)}\}$, and $s'\le t'$ (otherwise $s<t<s'\le s+2^{-(m+1)}$, which is
impossible in $D_{m+1}$). Since $t'-s'\le t-s<2^{-n}$, the induction hypothesis (or $t'=s'$) and
(c) at level $m+1>N$ give
\[|X_t-X_s|\le 2\sum_{j=n+1}^m2^{-\gamma j}+2\cdot2^{-\gamma(m+1)}.\]
Now let $s<t$ in $D=\bigcup_mD_m$ with $t-s<2^{-N(\omega)}$. Choose $n$ with
$2^{-(n+1)}\le t-s<2^{-n}$, so $n\ge N(\omega)$, and choose $m>n$ with $s,t\in D_m$. Then
\[
|X_t-X_s|\le 2\sum_{j>n}2^{-\gamma j}=\frac{2\cdot2^{-\gamma(n+1)}}{1-2^{-\gamma}}\le C_\gamma|t-s|^\gamma,
\qquad C_\gamma=\frac{2}{1-2^{-\gamma}}.
\]

\sub{(e)} Let $s<t$ in $D$ with $t-s\ge2^{-N}$. The points of $D_{N+1}\cap(s,t)$ together with $s,t$
form a chain $s=u_0<\dots<u_k=t$ in $D$ with consecutive gaps $<2^{-N}$ and $k\le2^{N+1}+1$.
Each gap contributes at most $C_\gamma2^{-N\gamma}\le C_\gamma|t-s|^\gamma$. Hence on $\{N<\infty\}$,
\[|X_t-X_s|\le K(\omega)|t-s|^\gamma\ \text{ for all } s,t\in D,\qquad K=(2^{N+1}+1)C_\gamma.\]

\sub{(f)} On $\Omega^*=\{N<\infty\}$, $X|_D$ is uniformly continuous, so
$Y_t:=\lim_{D\ni r\to t}X_r$ exists for all $t\in[0,1]$. Set $Y\equiv0$ off $\Omega^*$. Each $Y_t$ is
measurable (a limit along a sequence), and $Y$ is $\gamma$-H\"older on $[0,1]$ with constant $K$ by
continuity. For fixed $t$ take $r_j\in D$, $r_j\to t$: then $X_{r_j}\to Y_t$ a.s.\ and
$X_{r_j}\to X_t$ in probability by (b). Limits in probability are a.s.\ unique, so $Y_t=X_t$ a.s.,
and $Y$ is a modification of $X$. $Y$ does not depend on $\gamma$. Taking $\gamma_k\uparrow\beta/\alpha$
and intersecting the corresponding events $\Omega^*_{\gamma_k}$ shows that a.s.\ $Y$ is $\gamma$-H\"older for every
$\gamma<\beta/\alpha$.

\prob{2.}
Here $t_0$ is a local maximum if there is $\varepsilon>0$ with $B_s\le B_{t_0}$ for all
$s\in[0,\infty)\cap(t_0-\varepsilon,t_0+\varepsilon)$.

\sub{(a)} $W_s=B_{t_0+s}-B_{t_0}$ is a BM. If $t_0$ is a local maximum, then $\max_{[0,1/n]}W\le0$ for some $n$.
By the reflection principle,
\[\P\big(\max_{[0,1/n]}W\ge a\big)=2\P(W_{1/n}\ge a)\xrightarrow[a\downarrow0]{}1,\]
so $\P(\max_{[0,1/n]}W>0)=1$ for every $n$. Hence
$\P\big(\bigcup_n\{\max_{[0,1/n]}W\le0\}\big)=0$, and $\{t_0\text{ local max}\}$ is contained in a null set.

\sub{(b)} $B$ attains its maximum on $[0,1]$ at some $t^*$. By (a) with $t_0=0$, a.s.\ $t^*\neq0$. The process
$W_s=B_1-B_{1-s}$, $s\in[0,1]$, is a continuous centred Gaussian process. Its covariance is
\[
\mathrm{Cov}(W_s,W_u)=1-(1-u)-(1-s)+(1-s\vee u)=s\wedge u,
\]
so $W$ is a BM on $[0,1]$. Moreover $t^*=1$ means $\min_{[0,1]}W\ge0$, which has probability $0$ by the argument of (a)
applied to $-W$. So a.s.\ $t^*\in(0,1)$, and $t^*$ is a local maximum of $B$.

\prob{3.}
\sub{(a)} Let $Z_n=(B_{1/n}-B_{1/(n+1)})/\sqrt{1/n-1/(n+1)}$. The intervals $[\frac1{n+1},\frac1n]$ are
disjoint, so the $Z_n$ are i.i.d.\ $N(0,1)$. Then $\sum_n\P(Z_n\ge K)=\sum_n\P(Z_1\ge K)=\infty$, and the second Borel--Cantelli
lemma gives $\P(Z_n\ge K\text{ i.o.})=1$.

\sub{(b)} Fix rationals $p<q$. The process $W_t=B_{p+t}-B_p$ is a BM. Applying (a) to $W$ for every $K\in\N$, a.s.\ for all $K$
there are infinitely many $n$ with $1/n\le q-p$ and
\[|W_{1/n}-W_{1/(n+1)}|\ge K|1/n-1/(n+1)|^{1/2}.\]
So $B$ is not $\frac12$-H\"older on $[p,q]$. Intersecting over the countably many pairs $(p,q)$, and using that every
non-trivial interval contains some $[p,q]$, gives the claim.

\prob{4.}
Let $n\in\N$, $h=1/n$, $t_k=kh$, $\delta=\varepsilon/4$. Let $L$ be the piecewise linear interpolation of
$(t_k,B_{t_k})_k$ and $Z=B-L$. On $[t_{k-1},t_k]$,
\[
Z_s=(B_s-B_{t_{k-1}})-\tfrac{s-t_{k-1}}{h}(B_{t_k}-B_{t_{k-1}}),\qquad\text{so}\quad
\sup_{[t_{k-1},t_k]}|Z|\le2\sup_{u\in[0,h]}|B_{t_{k-1}+u}-B_{t_{k-1}}|.
\]
\emph{Independence.} $(Z_s)_s$ and $(B_{t_j})_j$ are jointly Gaussian and
$\mathrm{Cov}(Z_s,B_{t_j})=0$. For $t_j\le t_{k-1}$ this covariance equals
$t_j-\frac{(t_k-s)t_j+(s-t_{k-1})t_j}{h}=0$, and for $t_j\ge t_k$ it equals
$s-\frac{(t_k-s)t_{k-1}+(s-t_{k-1})t_k}{h}=0$. Hence $\sigma(Z_s:s\in[0,1])$ is independent of
$(B_{t_1},\dots,B_{t_n})$.

Let $f_h$ be the linear interpolation of $(t_k,f(t_k))$. Then $|f_h-f|\le\omega_f(h)$, where
$\omega_f$ is the modulus of continuity. Consider the events
\[
A=\{\max_k|B_{t_k}-f(t_k)|<\delta\},\qquad C=\{\sup_{[0,1]}|Z|<\delta\}.
\]
On $A\cap C$ we have $|B_t-f(t)|\le|L_t-f_h(t)|+|f_h(t)-f(t)|+|Z_t|<2\delta+\omega_f(h)$, which is
$<\varepsilon$ once $\omega_f(h)<\varepsilon/2$.

$\P(A)>0$, since $(B_{t_1},\dots,B_{t_n})$ is an invertible linear image of i.i.d.\ $N(0,h)$
increments and so has an everywhere positive density. The suprema over the different intervals are independent
(independent increments) and each has the law of $\sup_{[0,h]}|B|$. By scaling,
\[
\P(C)\ge\P\big(\sup_{[0,h]}|B|<\delta/2\big)^n=\P\big(\sup_{[0,1]}|B|<\delta/(2\sqrt h)\big)^n.
\]
The inner probability tends to $1$ as $h\to0$, so it is positive for $h$ small. Choosing such
$h$ with also $\omega_f(h)<\varepsilon/2$, independence gives
$\P(\sup_{[0,1]}|B-f|<\varepsilon)\ge\P(A)\P(C)>0$.

\prob{5.}
\sub{(a)} $X_0=0$ and $X$ has continuous paths. For $0\le t_0<\dots<t_k$, the increments
$X_{t_i}-X_{t_{i-1}}=a\cdot(B_{t_i}-B_{t_{i-1}})$ are functions of the independent vectors
$B_{t_i}-B_{t_{i-1}}$, hence independent. Each is a sum of independent
$N(0,a_j^2(t_i-t_{i-1}))$ variables, i.e.\ $N(0,|a|^2(t_i-t_{i-1}))=N(0,t_i-t_{i-1})$. So $X$ is a BM.

\sub{(b)} By (a) each coordinate $Y^i=a_i\cdot B$ is a BM, where $a_i$ is the $i$-th row. All $Y^i_t$ are linear combinations of the
Gaussian family $(B^j_t)$, so $(Y^1,\dots,Y^m)$ is jointly Gaussian with
$\mathrm{Cov}(Y^i_s,Y^j_t)=(a_i\cdot a_j)(s\wedge t)=\delta_{ij}(s\wedge t)$. Uncorrelated jointly Gaussian
families are independent: the finite-dimensional distributions factor, and a $\pi$--$\lambda$
argument finishes the proof. So $Y^1,\dots,Y^m$ are independent BMs.

\sub{(c)} For $R\in SO(d)$ the rows are orthonormal, so by (b) $RB$ consists of $d$ independent BMs, i.e.\
it is a $d$-dimensional BM. Since the law of a process is determined by its finite-dimensional distributions,
$(RB_t)_t\overset{d}{=}(B_t)_t$.

\prob{6.}
For a partition $\pi$ of $[0,T]$ write $Q_\pi(X)=\sum_i(X_{t_i}-X_{t_{i-1}})^2$. We use the known result
that $Q_\pi(B)\to b-a$ in $L^2$ on any interval $[a,b]$; indeed
$\E(Q_\pi(B)-(b-a))^2=2\sum\Delta t_i^2\le2(b-a)|\pi|$.

\sub{(a)} $Q_\pi(aB)=a^2Q_\pi(B)\to a^2T$.

\sub{(b)} $W_t=c^{-1/2}B_{ct}$ is a BM, so $Q_\pi(B_{c\cdot})=cQ_\pi(W)\to cT$.

\sub{(c)} Let $Y_t=\int_0^tB_s\,ds$ and $M=\sup_{[0,T]}|B|$. Then $|Y_{t_i}-Y_{t_{i-1}}|\le M\Delta t_i$, so
$Q_\pi(Y)\le M^2T|\pi|$ and $\E Q_\pi(Y)^2\le\E[M^4]T^2|\pi|^2\to0$. Here $\E M^4<\infty$, since
$M\le\sup B+\sup(-B)$ and $\sup_{[0,T]}B\overset d=|B_T|$ by the reflection principle. The quadratic variation is $0$.

\sub{(d)} If $t_0\in\pi$, then $Q_\pi(X)=Q_{\pi\cap[0,t_0]}(B)+9\,Q_{\pi\cap[t_0,T]}(B)\to t_0+9(T-t_0)$ in $L^2$.
Otherwise let $[t_{j-1},t_j]\ni t_0$, $u=B_{t_0}-B_{t_{j-1}}$, $v=-3(B_{t_j}-B_{t_0})$. Then
$Q_\pi(X)=Q_{\pi\cup\{t_0\}}(X)+2uv$, and by independence
$\E(2uv)^2=36(t_0-t_{j-1})(t_j-t_0)\le 36|\pi|^2\to0$. Hence the quadratic variation is $t_0+9(T-t_0)=9T-8t_0$.

%=====================================================================
\section*{Exercise 2}
\subsection*{Conditional expectation}

\prob{1.}
\sub{(a)} Let $Z=\E[X|Y]$. By conditional Jensen, $Z^2\le\E[X^2|Y]$, so $Z\in L^2$. If
$\E[h(Y)^2]=\infty$ the right-hand side is infinite (as $X\in L^2$). Otherwise $V=Z-h(Y)\in
L^2(\sigma(Y))$, and by Problem 2(b) (the product $V(X-Z)$ is integrable by Cauchy--Schwarz)
\[
\E[V(X-Z)]=\E\big[V\,\E[X-Z\,|\,Y]\big]=0 .
\]
Hence $\E(X-h(Y))^2=\E(X-Z)^2+\E V^2+2\E[V(X-Z)]\ge\E(X-Z)^2$.

\sub{(b)} $L^2(\Omega,\sigma(Y),\P)$ is a closed subspace of $L^2(\Omega,\F,\P)$. As shown, $Z=\E[X|Y]$ lies in
it, and $\E[V(X-Z)]=0$ for every $V\in L^2(\sigma(Y))$, i.e.\ $X-Z\perp L^2(\sigma(Y))$. These two
properties characterise the orthogonal projection. (Equivalently, by (a): every $\sigma(Y)$-measurable $V$ is of the form $h(Y)$, so $Z$ is the closest point.)

\prob{2.}
\sub{(a)} ($\Rightarrow$) Let $Z=\E[X|\G]$. For $W=\1_A$, $A\in\G$, the identity is the definition. By linearity
it holds for simple $W$. For $W$ with $|W|\le c$ a.s., take simple $\G$-measurable $W_n\to W$ with $|W_n|\le c$.
Then $|W_nZ|\le c|Z|$ and $|W_nX|\le c|X|$, and dominated convergence gives $\E[WZ]=\E[WX]$.

($\Leftarrow$) Taking $W=\1_A$ gives $\E[Z\1_A]=\E[X\1_A]$ for $A\in\G$. Integrability of $Z$: with
$W_n=\sgn(Z)\1_{\{|Z|\le n\}}$ we get $\E[|Z|\1_{\{|Z|\le n\}}]=\E[W_nX]\le\E|X|$. Letting $n\to\infty$ gives $\E|Z|\le\E|X|$ by
monotone convergence. Hence $Z=\E[X|\G]$.

\sub{(b)} First let $X,Y\ge0$ and $Z=\E[X|\G]\ge0$. For $A\in\G$, $W=\1_A\min(Y,n)$ is bounded and
$\G$-measurable, so $\E[\1_A\min(Y,n)Z]=\E[\1_A\min(Y,n)X]$. By monotone convergence,
$\E[\1_AYZ]=\E[\1_AYX]<\infty$. So $YZ$ is integrable, $\G$-measurable, and equals $\E[XY|\G]$.

In general write $X=X^+-X^-$ and $Y=Y^+-Y^-$. Each $X^\pm Y^\pm\le|XY|$ is integrable, and the
positive case gives $\E[X^\pm Y^\pm|\G]=Y^\pm\E[X^\pm|\G]$. Summing with signs and using linearity,
\[\E[XY|\G]=(Y^+-Y^-)(\E[X^+|\G]-\E[X^-|\G])=Y\,\E[X|\G].\]

\prob{3.}
$\G$ is generated by the partition $I_1=[0,\frac14]$, $I_2=(\frac14,\frac23]$, $I_3=(\frac23,1]$, so
$\G=\{\bigcup_{i\in S}I_i:S\subseteq\{1,2,3\}\}$. Let $Z=\sum_ia_i\1_{I_i}$ with
$a_i=\Leb(I_i)^{-1}\int_{I_i}\omega^2d\omega$. Then $Z$ is $\G$-measurable, and $\E[Z\1_{I_i}]=\E[X\1_{I_i}]$ for each $i$, hence for
every $A\in\G$ by additivity. By uniqueness $\E[X|\G]=Z$ a.s. The coefficients are
\[
a_1=4\cdot\tfrac{1}{3\cdot64}=\tfrac1{48},\qquad
a_2=\tfrac{12}{5}\cdot\tfrac13\Big(\tfrac{8}{27}-\tfrac1{64}\Big)=\tfrac{12}{5}\cdot\tfrac{485}{5184}=\tfrac{97}{432},\qquad
a_3=3\cdot\tfrac13\Big(1-\tfrac{8}{27}\Big)=\tfrac{19}{27}.
\]

\prob{4.}
\note{$Y$ must be allowed to take the value $+\infty$. For instance, if $\G=\{\emptyset,\Omega\}$ then
$\int_\Omega Y\,d\P=\E X=\infty$ forces $Y\equiv\infty$. We prove the statement for $[0,\infty]$-valued $Y$.}

\emph{Existence.} For $M\in\N$ let $X_M=\min(X,M)$ and $Y_M=\E[X_M|\G]$. Since $0\le X_M\le X_{M+1}$, monotonicity gives
$0\le Y_M\le Y_{M+1}$ a.s. Outside a null set the sequence is non-decreasing, and $Y:=\lim_MY_M\in[0,\infty]$ is
$\G$-measurable. For $A\in\G$, monotone convergence gives
\[\int_AY=\lim_M\int_AY_M=\lim_M\int_AX_M=\int_AX.\]

\emph{Uniqueness.} Let $Y,Y'$ be two such variables and $A_n=\{Y'>Y,\ Y\le n\}\in\G$. Then
$\int_{A_n}Y'=\int_{A_n}X=\int_{A_n}Y\le n<\infty$, so $\int_{A_n}(Y'-Y)=0$ with $Y'-Y>0$ on $A_n$. Hence
$\P(A_n)=0$, and taking the union over $n$ gives $\P(Y'>Y,\,Y<\infty)=0$. Since $Y'>Y$ forces $Y<\infty$, we get $\P(Y'>Y)=0$.
By symmetry $\P(Y>Y')=0$.

\prob{5.}
Take $\Omega=\{a,b,c\}$, $\F=2^\Omega$, $\P$ uniform, $X=\1_{\{a\}}$, $\G=\sigma(\{a\})$, $\Hc=\sigma(\{c\})$. Then
$\E[X|\G]=X$ and
\[
\E[\E[X|\G]|\Hc]=\tfrac12\1_{\{a,b\}},\qquad
\E[X|\Hc]=\tfrac12\1_{\{a,b\}},\qquad \E[\E[X|\Hc]|\G]=\tfrac12\1_{\{a\}}+\tfrac14\1_{\{b,c\}}.
\]
At $c$ (which has probability $\frac13$) the first equals $0$ and the second equals $\frac14$.

\subsection*{Martingales}

\prob{1.}
($\Rightarrow$) For $A\in\F_s$, $\E[X_t\1_A]=\E[\E[X_t|\F_s]\1_A]\ge\E[X_s\1_A]$.
($\Leftarrow$) Let $Z=\E[X_t|\F_s]-X_s$, which is $\F_s$-measurable with $\E[Z\1_A]\ge0$ for all $A\in\F_s$.
With $A=\{Z<0\}$ we get $\E[Z\1_{\{Z<0\}}]\ge0$, so $\P(Z<0)=0$, i.e.\ $\E[X_t|\F_s]\ge X_s$.

\prob{2.}
(The index set in $(X_t,\F^X_t)_{t\in X}$ should be $T$.) ($\Leftarrow$) Take $\F_t=\F^X_t$.
($\Rightarrow$) Since $X$ is $(\F_t)$-adapted, $\F^X_t\subseteq\F_t$, and $X$ is $\F^X$-adapted and integrable.
For $s<t$, by the tower property,
\[\E[X_t|\F^X_s]=\E[\E[X_t|\F_s]|\F^X_s]=\E[X_s|\F^X_s]=X_s.\]

\prob{3.}
(Doob's theorem; (a) should read $\E|X_{\tau_2}|<\infty$.) Here $\F_\infty\supseteq\F_n$ for all $n$, and the martingale
property gives $X_n=\E[X_\infty|\F_n]$. A stopping time is $\tau:\Omega\to\N_\infty$ with $\{\tau=n\}\in\F_n$, $n\in\N$, and
$\F_\tau=\{A\in\F_\infty:A\cap\{\tau=n\}\in\F_n\ \forall n\in\N\}$.

\emph{Claim:} for every stopping time $\tau$, $X_\tau=\E[X_\infty|\F_\tau]$.

$X_\tau=\sum_{n\in\N_\infty}X_n\1_{\{\tau=n\}}$ is $\F_\tau$-measurable, since
$\{X_\tau\in B\}\cap\{\tau=n\}=\{X_n\in B\}\cap\{\tau=n\}\in\F_n$. By conditional Jensen, $|X_n|\le\E[|X_\infty||\F_n]$, so
\[\E|X_\tau|\le\sum_n\E[|X_\infty|\1_{\{\tau=n\}}]=\E|X_\infty|.\]
For $A\in\F_\tau$, Fubini (justified by the last bound) gives
\[\E[X_\tau\1_A]=\sum_n\E[X_n\1_{A\cap\{\tau=n\}}]=\sum_n\E[X_\infty\1_{A\cap\{\tau=n\}}]=\E[X_\infty\1_A],\]
because $A\cap\{\tau=n\}\in\F_n$ (trivially for $n=\infty$).

\sub{(a)} $\E|X_{\tau_2}|\le\E|X_\infty|<\infty$.

\sub{(b)} If $A\in\F_{\tau_1}$ then $A\cap\{\tau_2\le n\}=A\cap\{\tau_1\le n\}\cap\{\tau_2\le n\}\in\F_n$, so
$\F_{\tau_1}\subseteq\F_{\tau_2}$. By the claim and the tower property,
\[\E[X_{\tau_2}|\F_{\tau_1}]=\E[\E[X_\infty|\F_{\tau_2}]|\F_{\tau_1}]=\E[X_\infty|\F_{\tau_1}]=X_{\tau_1}.\]

\prob{4.}
$\F_\tau=\{A:A\cap\{\tau\le u\}\in\F_u\ \forall u\}$. From discrete time we use: if $(Z_g,\F_g)_{g\in G}$ is a
supermartingale on a discrete set $G\subset[0,\infty)$ and $\sigma\le\rho\le c$ are $G$-valued stopping
times, then $\E[Z_\rho|\F_\sigma]\le Z_\sigma$, and the stopped process $(Z_{g\wedge\rho})_{g\in G}$ is a
supermartingale.

\emph{Lemma.} Let $Y_n$ be integrable and $\Hc_n$-measurable with $\E[Y_m\1_B]\le\E[Y_n\1_B]$ for all $m\le n$ and
$B\in\Hc_n$, and suppose $L=\sup_n\E Y_n<\infty$. Then $(Y_n)$ is uniformly integrable.

\emph{Proof.} Taking $B=\Omega$ shows $\E Y_n\uparrow L$. From $Y_n\ge\E[Y_1|\Hc_n]\ge-\E[Y_1^-|\Hc_n]$ we get $\E Y_n^-\le\E Y_1^-$, so
$\sup_n\E|Y_n|\le L+2\E|Y_1|=:K$. Given $\varepsilon>0$ pick $m$ with $\E Y_m\ge L-\varepsilon$. For $n\ge m$, since
$\{Y_n\le\lambda\},\{Y_n<-\lambda\}\in\Hc_n$,
\begin{align*}
\E[|Y_n|;|Y_n|>\lambda]&=\E Y_n-\E[Y_n;Y_n\le\lambda]-\E[Y_n;Y_n<-\lambda]\\
&\le\E Y_m+\varepsilon-\E[Y_m;Y_n\le\lambda]-\E[Y_m;Y_n<-\lambda]\le\E[|Y_m|;|Y_n|>\lambda]+\varepsilon .
\end{align*}
As $\P(|Y_n|>\lambda)\le K/\lambda$ and $Y_m$ is integrable, the right side is $\le2\varepsilon$ for $\lambda$ large, uniformly in $n\ge m$.
The finitely many $n<m$ are handled trivially. \hfill$\square$

\sub{(a)} Fix $s<t$ and let $G_n=2^{-n}\N\cup\{s,t\}$; these sets are nested. Let $\tau_n=\inf\{g\in G_n:g>\tau\}$, with $\tau_n=\infty$ if
$\tau=\infty$. Then $\{\tau_n\le g\}=\{\tau<g\}\in\F_g$, so $\tau_n$ is a $G_n$-valued stopping time, and
it is also a stopping time for the grid $G_{n'}$, $n'\ge n$. Moreover $\tau<\tau_n\le\tau+2^{-n}$ and $\tau_n\downarrow\tau$.

For $r\in\{s,t\}$, $X_{r\wedge\tau_n}=\sum_{g\in G_n\cap[0,r]}X_g\1_{\{r\wedge\tau_n=g\}}$ is
$\F_r$-measurable. By right-continuity $X_{r\wedge\tau_n}\to X_{r\wedge\tau}$, so $X_{r\wedge\tau}$ is $\F_r$-measurable.
For $u<r$, $\{X_{r\wedge\tau}\in B\}\cap\{\tau\le u\}=\{X_{u\wedge\tau}\in B\}\cap\{\tau\le u\}\in\F_u$, hence
$X_{r\wedge\tau}$ is $\F_{r\wedge\tau}$-measurable.

Since the stopped discrete process is a supermartingale and $s,t\in G_n$, for $A\in\F_s$,
\[
\E[X_{t\wedge\tau_n}\1_A]\le\E[X_{s\wedge\tau_n}\1_A].\tag{$*$}
\]
Let $Y_n=X_{r\wedge\tau_n}$ and let $\Hc_n$ be the $\sigma$-algebra of $r\wedge\tau_n$ on the grid $G_n$. For $m\le n$, $r\wedge\tau_n\le
r\wedge\tau_m\le r$ are $G_n$-stopping times, so $\E[Y_m|\Hc_n]\le Y_n$. Also $\E Y_n\le\E X_0$. By the Lemma,
$(Y_n)$ is UI. Together with a.s.\ convergence this gives $X_{r\wedge\tau_n}\to X_{r\wedge\tau}$ in $L^1$, so
$X_{r\wedge\tau}$ is integrable. Passing to the limit in $(*)$ gives
$\E[X_{t\wedge\tau}\1_A]\le\E[X_{s\wedge\tau}\1_A]$ for $A\in\F_s$. By Problem 1 (with reversed inequality),
$(X_{t\wedge\tau},\F_t)$ is a supermartingale.

Since $\F_{s\wedge\tau}\subseteq\F_s$ (if $A\in\F_{s\wedge\tau}$ then $A=A\cap\{s\wedge\tau\le s\}\in\F_s$),
$(\F_{t\wedge\tau})$ is a filtration to which $X_{t\wedge\tau}$ is adapted. The same inequality for
$A\in\F_{s\wedge\tau}$ shows that $(X_{t\wedge\tau},\F_{t\wedge\tau})$ is a supermartingale.

\sub{(b)} Applying (a) to $X$ and $-X$, $(X_{t\wedge\tau})$ is a martingale, so $\E X_{t\wedge\tau}=\E X_0$.
On $\{\tau<\infty\}$, $X_{t\wedge\tau}=X_\tau$ for $t\ge\tau$. On $\{\tau=\infty\}$, $X_\tau$ must be given a meaning. We
set $X_\tau:=\lim_tX_{t\wedge\tau}$, which exists a.s.\ by Problem 5(a) because $\sup_t\E X_{t\wedge\tau}^+\le\E Y$. (If
$\tau<\infty$ a.s.\ this is not needed.) Then $|X_\tau|\le Y$, so $X_\tau$ is integrable, and dominated convergence gives
$\E X_\tau=\lim_t\E X_{t\wedge\tau}=\E X_0$.

\prob{5.}
\sub{(a)} Let $D$ be the non-negative dyadic rationals and $F_k=D_k\cap[0,k]$. For $a<b$ let
$U_k$ be the number of upcrossings of $[a,b]$ by $(X_t)_{t\in F_k}$. By Doob's upcrossing inequality for the
discrete submartingale on the finite set $F_k$,
\[(b-a)\E U_k\le\E(X_k-a)^+\le\sup_t\E X_t^++|a|.\]
$U_k$ increases to the number $U_D$ of upcrossings by $(X_t)_{t\in D}$, so by monotone convergence $\E U_D<\infty$. Hence a.s.\
$U_D[a,b]<\infty$ for all rational $a<b$, and $\lim_{D\ni t\to\infty}X_t$ exists in $[-\infty,\infty]$ (otherwise
$\liminf<a<b<\limsup$ would give infinitely many upcrossings).

By right-continuity, if $X_t>c$ for arbitrarily large real $t$ then $X_d>c$ for arbitrarily large $d\in D$. So
$\limsup_{t\to\infty}X_t=\limsup_{D\ni t\to\infty}X_t$, and similarly for $\liminf$. Thus $X=\lim_{t\to\infty}X_t$ exists a.s.
Since $\E|X_t|=2\E X_t^+-\E X_t\le2\sup_t\E X_t^+-\E X_0$, Fatou's lemma gives $\E|X|<\infty$.

\sub{(b)} \note{Right-continuity must be assumed here as well. Let $U\sim U[0,1]$ and
$X_t=\1_{\{t-\lfloor t\rfloor=U\}}$. Then $X_t=0$ a.s.\ for each $t$, so $X$ is an $L^2$-bounded martingale, but
every path equals $1$ at $U+k$ for all $k\in\N$ and does not converge.}

Assume right-continuity. Since $\sup\E X_t^+\le\sup(\E X_t^2)^{1/2}<\infty$, (a) gives $X_t\to X$ a.s.
For $s<t$, by Conditional expectation, Problem 2(b), $\E[X_s(X_t-X_s)]=\E[X_s\E[X_t-X_s|\F_s]]=0$. Hence
$\E(X_t-X_s)^2=\E X_t^2-\E X_s^2$. So $\E X_t^2$ is non-decreasing and bounded, and $\E(X_t-X_s)^2\to0$ as $s,t\to\infty$.
By completeness of $L^2$, $X_t\to X'$ in $L^2$. Along a sequence $t_n\to\infty$ a subsequence converges a.s.\ to $X'$,
so $X'=X$ a.s.\ and $X\in L^2$. Finally, for $u>t$, $X_t=\E[X_u|\F_t]$ and
\[\|\E[X_u|\F_t]-\E[X|\F_t]\|_1\le\|X_u-X\|_1\to0,\]
so $\E[X|\F_t]=X_t$ a.s.

\prob{6.}
\sub{(a)} The argument works for every $\sigma\in\R$, which we need in (b). Let $M_t(\sigma)=\exp(\sigma
B_t-\sigma^2t/2)$. It is adapted, and integrable since $\E e^{\sigma B_t}=e^{\sigma^2t/2}$. Since $B_t-B_s$ is independent
of $\F_s$,
\[\E[M_t|\F_s]=M_s\,\E\big[e^{\sigma(B_t-B_s)}\big]e^{-\sigma^2(t-s)/2}=M_s.\]

\sub{(b)} Writing $x=B_t$,
$e^{\sigma x-\sigma^2t/2}=e^{\sigma x}e^{-\sigma^2t/2}=\sum_k\frac{\sigma^k}{k!}H_k(x,t)$, where
\[
H_2=x^2-t,\quad H_3=x^3-3tx,\quad H_4=x^4-6tx^2+3t^2
\]
are read off from the coefficients of $\sigma^2,\sigma^3,\sigma^4$ in
$(1+\sigma x+\frac{\sigma^2x^2}2+\frac{\sigma^3x^3}6+\frac{\sigma^4x^4}{24})(1-\frac{\sigma^2t}2+\frac{\sigma^4t^2}8)$.
Thus $\partial_\sigma^kM_t(\sigma)|_{\sigma=0}=H_k(B_t,t)$.

By (a), $\E[M_t(\sigma)\1_A]=\E[M_s(\sigma)\1_A]$ for $A\in\F_s$ and all $\sigma$. For $|\sigma|\le1$,
$\partial^j_\sigma M_t(\sigma)=P_j(B_t-\sigma t,t)M_t(\sigma)$ with $P_j$ a polynomial. This is dominated by
$C_j(1+|B_t|+t)^je^{|B_t|}$, which is integrable, so we may differentiate under $\E$ $k$ times at $\sigma=0$:
$\E[H_k(B_t,t)\1_A]=\E[H_k(B_s,s)\1_A]$. Since $H_k(B_t,t)$ is adapted and integrable, it is a martingale. For
$k=2,3,4$ this gives the processes $B_t^2-t$, $B_t^3-3tB_t$ and $B_t^4-6tB_t^2+3t^2$.

%=====================================================================
\section*{Exercise 4}
\subsection*{Stochastic calculus}

Throughout, for $f\in L^2$ (adapted, with $\E\int_0^tf^2ds<\infty$), $\int_0^tf\,dB$ is the $L^2$-limit of
$\int f_n\,dB$ for simple adapted $f_n\to f$ in $L^2(ds\otimes d\P)$. For a partition
$0=t_0<\dots<t_n=t$ of mesh $|\pi|$ write $\Delta_i=B_{t_i}-B_{t_{i-1}}$ and $\Delta t_i=t_i-t_{i-1}$.

\prob{1.}
\note{In (a)ii and (b)i the right-hand side should be $\frac13B_t^3-\int_0^tB_s\,ds$, not $\frac13B_t^2-\dots$.
As stated it cannot hold: the left side has mean $0$, while $\E[\frac13B_t^2-\int_0^tB_s\,ds]=\frac t3$.
Itô's formula in (b) confirms the corrected form.}

\sub{(a) i.} $f_\pi=\sum t_{i-1}\1_{[t_{i-1},t_i)}\to s$ in $L^2$, since $\int_0^t|f_\pi(s)-s|^2ds\le t|\pi|^2$. Summation by parts gives
\[
\sum_it_{i-1}\Delta_i=tB_t-\sum_iB_{t_i}\Delta t_i .
\]
By continuity $|\sum_iB_{t_i}\Delta t_i-\int_0^tB_s\,ds|\to0$ pathwise. It is bounded by $2t\sup_{[0,t]}|B|\in L^2$,
so the convergence also holds in $L^2$. Hence $\int_0^ts\,dB_s=tB_t-\int_0^tB_s\,ds$.

\sub{(a) ii.} For $u<s\le t$, by Cauchy--Schwarz and $B_s+B_u\sim N(0,s+3u)$,
$\E(B_s^2-B_u^2)^2\le\sqrt{3(s-u)^2}\sqrt{3(4t)^2}=12t(s-u)$. Hence
$\sum_iB_{t_{i-1}}^2\1_{[t_{i-1},t_i)}\to B^2$ in $L^2(ds\otimes d\P)$, and $\int_0^tB^2dB=\lim_{L^2}S_\pi$
with $S_\pi=\sum_iB_{t_{i-1}}^2\Delta_i$. Summing $b^3-a^3=3a^2(b-a)+3a(b-a)^2+(b-a)^3$ over the partition gives
\[
S_\pi=\tfrac13B_t^3-\sum_iB_{t_{i-1}}\Delta_i^2-\tfrac13\sum_i\Delta_i^3 .
\]
The cross terms below vanish by independence of increments and since odd moments are $0$, resp.\ $\E[\Delta_j^2-\Delta t_j]=0$:
\[
\E\Big(\sum_i\Delta_i^3\Big)^2=\sum_i15\Delta t_i^3\le15t|\pi|^2,\qquad
\E\Big(\sum_iB_{t_{i-1}}(\Delta_i^2-\Delta t_i)\Big)^2=\sum_it_{i-1}\cdot2\Delta t_i^2\le2t^2|\pi| .
\]
Moreover $\sum_iB_{t_{i-1}}\Delta t_i\to\int_0^tB_s\,ds$ in $L^2$ as in i. Hence
$\int_0^tB_s^2\,dB_s=\frac13B_t^3-\int_0^tB_s\,ds$.

\sub{(b) i.} Itô's formula with $f(x)=x^3/3$ (so $f'=x^2$, $f''=2x$) gives
$\frac13B_t^3=\int_0^tB_s^2\,dB_s+\int_0^tB_s\,ds$.

\prob{2.}
\note{We take $f,g\in L^2([0,T]\times\Omega)$, i.e.\ $\E\int_0^Tf^2<\infty$ for every $T$, the class on which the Itô integral is
an $L^2$-martingale. If only $\int_0^Tf^2ds<\infty$ a.s.\ were assumed, the claim is false in
general; we do not pursue this.}

The Itô integral has mean zero, so taking expectations gives $C=D$. Then $\int_0^T(f-g)\,dB=0$ a.s., and by the Itô
isometry
\[\E\int_0^T(f-g)^2ds=\E\Big(\int_0^T(f-g)\,dB\Big)^2=0,\]
i.e.\ $f=g$ for $ds\otimes d\P$-a.e.\ $(s,\omega)$.

\prob{3.}
\sub{(a)} $d(B_t^2)=dt+2B_t\,dB_t$.

\sub{(b)} $d(2+t+e^{B_t})=(1+\tfrac12e^{B_t})\,dt+e^{B_t}\,dB_t$.

\sub{(c)} $d(B_t^2+\tilde B_t^2)=2\,dt+2B_t\,dB_t+2\tilde B_t\,d\tilde B_t$, i.e.\ $\alpha=2$ and $\beta=(2B_t,2\tilde B_t)$ against
$d(B_t,\tilde B_t)$.

\sub{(d)} $d(t,B_t)=\binom10dt+\binom01dB_t$.

\sub{(e)} With $f(t,x)=e^{t/2}\cos x$: $f_t=\frac12f$, $f_x=-e^{t/2}\sin x$, $f_{xx}=-f$. Hence
\[
d\big(e^{t/2}\cos B_t\big)=\big(\tfrac12f-\tfrac12f\big)dt-e^{t/2}\sin B_t\,dB_t=-e^{t/2}\sin B_t\,dB_t .
\]
So $e^{t/2}\cos B_t=1-\int_0^te^{s/2}\sin B_s\,dB_s$. The integrand is bounded by $e^{T/2}$ on $[0,T]$, so it lies in $L^2$ and
the Itô integral is a martingale. Therefore $e^{t/2}\cos B_t$ is a martingale.

\prob{4.}
\sub{(a)} With $f=e^x$, $f'=f''=f$: $dX_t=e^{B_t}dB_t+\frac12e^{B_t}dt=\frac12X_tdt+X_tdB_t$.

\sub{(b)} With $f=\sin x$: $dX_t=\cos B_t\,dB_t-\frac12\sin B_t\,dt$. While $\cos B_t>0$, i.e.\ for
$t<\inf\{s:|B_s|=\pi/2\}$, we have $\cos B_t=\sqrt{1-\sin^2B_t}=\sqrt{1-X_t^2}$. Hence
$dX_t=-\frac12X_tdt+\sqrt{1-X_t^2}\,dB_t$.

\prob{5.}
For $k\ge2$, Itô's formula gives
\[B_t^k=\int_0^tkB_s^{k-1}dB_s+\tfrac12k(k-1)\int_0^tB_s^{k-2}ds.\]
The integrand satisfies $\E\int_0^tB_s^{2k-2}ds<\infty$, so the stochastic integral has mean $0$. By Fubini,
\[\beta_k(t)=\tfrac12k(k-1)\int_0^t\beta_{k-2}(s)\,ds.\]
The formula also holds for $k=1$ ($\beta_1=0$) but fails for $k=0$, where $\beta_0=1$. With $\beta_0=1$:
$\beta_2(t)=t$, then $\E B_t^4=\beta_4(t)=6\int_0^ts\,ds=3t^2$ and $\E B_t^6=\beta_6(t)=15\int_0^t3s^2\,ds=15t^3$.

\prob{6.}
\note{$g_\varepsilon$ is not $C^2$ on $\R$. It is $C^1$, with
$g_\varepsilon'(x)=\sgn x$ for $|x|\ge\varepsilon$ and $x/\varepsilon$ for $|x|<\varepsilon$. Its second derivative
$g_\varepsilon''=\varepsilon^{-1}\1_{(-\varepsilon,\varepsilon)}$ exists except at $\pm\varepsilon$, so $g_\varepsilon$ is
$C^2$ on $\R\setminus\{\pm\varepsilon\}$. Itô's formula still applies, as shown in (a).}

\sub{(a)} Let $\phi_n$ be continuous, $0\le\phi_n\le1/\varepsilon$, with $\phi_n=g_\varepsilon''$ outside
$(\pm\varepsilon-\frac1n,\pm\varepsilon+\frac1n)$. Let $h_n(x)=g_\varepsilon(0)+\int_0^x\int_0^y\phi_n$, which is $C^2$. Since
$g_\varepsilon'(0)=0$, $|h_n'-g_\varepsilon'|\le\int|\phi_n-g_\varepsilon''|\le4/(n\varepsilon)$ and
$|h_n-g_\varepsilon|(x)\le4|x|/(n\varepsilon)$. Itô's formula for $h_n$ reads
\[
h_n(B_t)=h_n(B_0)+\int_0^th_n'(B_s)\,dB_s+\tfrac12\int_0^t\phi_n(B_s)\,ds .
\]
As $n\to\infty$: $h_n(B_t)\to g_\varepsilon(B_t)$, and $\int h_n'(B)dB\to\int g_\varepsilon'(B)dB$ in $L^2$ by the isometry and
uniform convergence. Also $\phi_n(B_s)\to g_\varepsilon''(B_s)$ for all $s$ with $B_s\ne\pm\varepsilon$, which is a.e.\ $s$ a.s.\
because $\E\,\Leb\{s\le t:B_s=\pm\varepsilon\}=\int_0^t\P(B_s=\pm\varepsilon)ds=0$. Bounded convergence then gives
\[
g_\varepsilon(B_t)=g_\varepsilon(B_0)+\int_0^tg_\varepsilon'(B_s)\,dB_s+\frac1{2\varepsilon}\Leb\{s\in[0,t]:B_s\in(-\varepsilon,\varepsilon)\}.
\]

\sub{(b)} By the isometry and dominated convergence ($\P(B_s=0)=0$ for $s>0$),
\[
\E\Big(\int_0^t\tfrac{B_s}{\varepsilon}\1_{\{|B_s|<\varepsilon\}}dB_s\Big)^2=\E\int_0^t\tfrac{B_s^2}{\varepsilon^2}\1_{\{|B_s|<\varepsilon\}}ds
\le\int_0^t\P(|B_s|<\varepsilon)\,ds\xrightarrow[\varepsilon\to0]{}0 .
\]

\sub{(c)} Split $\int g_\varepsilon'(B)dB=\int\sgn(B)\1_{\{|B|\ge\varepsilon\}}dB+\int\frac{B}{\varepsilon}\1_{\{|B|<\varepsilon\}}dB$.
The second term tends to $0$ in $L^2$ by (b). The first tends to $\int_0^t\sgn(B_s)dB_s$ in $L^2$, since
$\E\int_0^t\1_{\{|B_s|<\varepsilon\}}ds\to0$. Moreover $0\le g_\varepsilon(x)-|x|=(\varepsilon-|x|)^2/(2\varepsilon)\,\1_{\{|x|<\varepsilon\}}\le\varepsilon/2$.
Hence by (a)
\[
\frac1{2\varepsilon}\Leb\{s\le t:B_s\in(-\varepsilon,\varepsilon)\}=g_\varepsilon(B_t)-g_\varepsilon(B_0)-\int_0^tg_\varepsilon'(B_s)dB_s
\]
converges in $L^2$ to $|B_t|-|B_0|-\int_0^t\sgn(B_s)dB_s$. Thus $L_t$ exists as an $L^2$-limit and
$|B_t|=|B_0|+\int_0^t\sgn(B_s)dB_s+L_t$.

\prob{7.}
\note{The definition should read $T_x=\inf\{t>0:X^x_t=0\}$ with $x>0$, not $\inf\{t:X^x_t>0\}$. In (b) the claim is
$(Y_t)\overset d=(X^1_t)$. All $X^x$, $x>0$, are driven by the same $B$.}

Facts used: $X^x$ is continuous and positive on $[0,T_x)$; $B_t/t\to0$ a.s.; and the
following consequence of the equation. If $X^x_t<R$ for all $t\ge0$, then $T_x=\infty$ and
$X^x_t\ge x+B_t+at/R\to\infty$, a contradiction. So $X^x$ cannot stay in $(r,R)$ forever.

\sub{(a)} For $t<S=T_x\wedge T_y$, $D_t=X^y_t-X^x_t=y-x+a\int_0^t\big(\frac1{X^y}-\frac1{X^x}\big)ds$ is $C^1$
with $D'=-aD/(X^xX^y)$. Hence
\[D_t=(y-x)\exp\Big(-a\int_0^t\frac{ds}{X^x_sX^y_s}\Big)>0 .\]
If $T_y<T_x$, then $0<X^x_t<X^y_t\to0$ as $t\uparrow T_y$, so $X^x_{T_y}=0$, a contradiction. So $T_x\le T_y$ and
$X^x_t<X^y_t$ for $t<T_x$. With $y=x$, the same computation gives $D\equiv0$: the equation has at most one
solution for a given path of $B$ (pathwise uniqueness).

\sub{(b)} Let $W_t=x^{-1}B_{x^2t}$, a BM. Substituting $s=x^2u$ in
$\int_0^{x^2t}ds/X^x_s=x^2\int_0^tdu/X^x_{x^2u}=x\int_0^tdu/Y_u$,
\[
Y_t=1+W_t+a\int_0^t\frac{du}{Y_u},\qquad t<T_x/x^2 .
\]
So $Y$ solves the equation with $x=1$ driven by $W$, and its lifetime is $T_x/x^2$. By pathwise uniqueness the solution is
a fixed functional of the driving path, $X^1=\Phi(B)$, and therefore $Y=\Phi(W)\overset d=\Phi(B)=X^1$.

\sub{(c)} Let $s(x)=x^{1-2a}$ for $a\ne\frac12$ and $s(x)=\log x$ for $a=\frac12$. Then $\frac12s''+\frac ax s'=0$. For
$0<r<x<R$ let $\tau$ be the exit time of $X^x$ from $(r,R)$; it is a.s.\ finite by the facts above. Itô's
formula for $dX=\frac a Xdt+dB$ gives $s(X_{t\wedge\tau})=s(x)+\int_0^{t\wedge\tau}s'(X)dB$, whose integrand is bounded,
so the integral is a martingale. Bounded convergence yields $\E s(X_\tau)=s(x)$, i.e.\
\[
p_{r,R}(x):=\P(X^x\text{ hits }r\text{ before }R)=\frac{s(R)-s(x)}{s(R)-s(r)} .
\]
By (a), $T_x\le T_y$ and $X^x\le X^y$ for $x<y$. Hence it suffices to prove i, ii and iv for the countably many $x\in\N$,
and iii for $x=1/n$, $n\in\N$.

\emph{i.} ($a\ge\frac12$.) If $T_x<\infty$, then $\sup_{[0,T_x]}X^x<R$ for some $R$, so $X^x$ hits every $r$ before $R$. But
$p_{r,R}(x)\to0$ as $r\to0$, since $s(r)\to-\infty$ for $a=\frac12$, and for $a>\frac12$ we have $s(r)\to+\infty$ while
$s(R)-s(x)$ is finite. Hence $\P(T_x<\infty)\le\sum_R\lim_{r\to0}p_{r,R}(x)=0$.

\emph{ii.} ($a=\frac12$.) By i, $T_x=\infty$. $\P(X^x\text{ hits }r)\ge p_{r,R}(x)=\frac{\log R-\log x}{\log R-\log r}\to1$ as
$R\to\infty$. Intersecting over $r=1/n$ gives $\inf_tX^x_t=0$ a.s.

\emph{iii.} ($a>\frac12$.) Here $T_x=\infty$ and $M_t=(X^x_t)^{1-2a}>0$. With
$\sigma_n=\inf\{t:X^x_t\notin(\frac1n,n)\}\uparrow\infty$, the process $M_{t\wedge\sigma_n}$ is a bounded martingale. Conditional
Fatou gives $\E[M_t|\F_s]\le\liminf_n M_{s\wedge\sigma_n}=M_s$, so $M$ is a non-negative continuous supermartingale. By
Exercise 2, Martingales 5(a) applied to $-M$, $M_t\to M_\infty\in[0,\infty)$ a.s., hence $X^x_t\to X_\infty\in(0,\infty]$.
If $X_\infty<\infty$ then $X^x$ is bounded, which contradicts the facts above. So $X^x_t\to\infty$.

\emph{iv.} ($a<\frac12$.) Here $s(0)=0$ and $s(\infty)=\infty$. The events $E_r=\{X^x\text{ hits }r\text{ before }R\}$ decrease as $r\downarrow0$.
On $\bigcap_rE_r$ the hitting times $\tau_r$ increase to some $\sigma\le\tau_R$. Moreover $\sigma<\infty$ a.s.: $\tau_R<\infty$ when $T_x=\infty$ by
the facts above, and $\sigma\le T_x$ otherwise. By continuity $X^x_{\sigma}=0$, so $\sigma=T_x<\tau_R$. Hence
\[\P(T_x<\tau_R)=\lim_{r\to0}p_{r,R}(x)=1-(x/R)^{1-2a}\to1\quad(R\to\infty),\]
i.e.\ $T_x<\infty$ a.s.

For v and vi fix $y>x>0$, write $X=X^x$, $D=X^y-X^x$, and let $V=X/D$ on $[0,T_x)$. $D$ is $C^1$, $D>0$ on $[0,T_x)$, and
$T_x<\infty$ a.s.\ by iv. By (a), $T_x=T_y$ iff $D_{T_x}=0$. From $d(1/D)=\frac{a}{XX^yD}dt$, $X=DV$ and $X^y=D(V+1)$,
\[
dV=\frac1D\Big(\frac aXdt+dB\Big)+X\,d\Big(\frac1D\Big)
=\frac{1}{D^2}\Big(\frac aV+\frac a{V+1}\Big)dt+\frac1D\,dB .
\]
For $\varphi\in C^2(0,\infty)$, Itô's formula gives
$d\varphi(V)=D^{-2}\big(\tfrac12\varphi''+(\tfrac aV+\tfrac a{V+1})\varphi'\big)(V)\,dt+\varphi'(V)D^{-1}dB$. Let
\[
S(v)=\int_0^v\big(u(1+u)\big)^{-2a}du,
\]
which is finite since $a<\frac12$. Then $\frac12S''+(\frac av+\frac a{v+1})S'=0$, and $S(V_t)$ has no $dt$ term.
Let $\rho_n=\inf\{t<T_x:D_t\le\frac1n\text{ or }V_t\le\frac1n\}\wedge n\wedge T_x$. The integrand of $S(V)$ is bounded up
to $\rho_n$, so $\E S(V_{\rho_n\wedge\kappa})=S(V_0)$ for every stopping time $\kappa$. Moreover $\rho_n\uparrow T_x$:
if the limit were $<T_x$, then $D$ and $V$ would be bounded away from $0$ on a compact interval past it.

\emph{v.} ($\frac14<a<\frac12$.) Now $S(\infty)<\infty$, since $(u(1+u))^{-2a}\sim u^{-4a}$ and $4a>1$. Suppose $\P(T_x<T_y)=1$.
On $\{T_x<T_y\}$ we have $D_{T_x}>0$ and $X_t\to0$, so $V_t\to0$ as $t\uparrow T_x$. Hence for large $n$, $\rho_n$ is the hitting
time of $\frac1n$ by $V$ and $S(V_{\rho_n})=S(\frac1n)\to0$. Since $0\le S\le S(\infty)$, bounded convergence gives
$S(V_0)=\lim_n\E S(V_{\rho_n})=0$, contradicting $V_0=x/(y-x)>0$. Hence $\P(T_x=T_y)>0$.

\emph{vi.} ($a\le\frac14$.) Now $S(\infty)=\infty$. Let $\kappa_\lambda=\inf\{t:S(V_t)\ge\lambda\}$. Since $S\ge0$,
$\lambda\P(\kappa_\lambda\le\rho_n)\le\E S(V_{\rho_n\wedge\kappa_\lambda})=S(V_0)$. Letting $n\to\infty$ and then $\lambda\to\infty$ gives
$\sup_{t<T_x}V_t<\infty$ a.s. Now fix $K$ and let $\rho_n'=\inf\{t:D_t\le\frac1n\text{ or }V_t\ge K\}\wedge n\wedge T_x$. Taking
expectations in the equation for $V$ (with $d\sigma=dt/D^2$, and the stochastic integral having bounded integrand $1/D\le n$),
\[
\E\int_0^{\rho_n'}\frac{a}{V_t}\frac{dt}{D_t^2}\le\E V_{\rho_n'}\le K .
\]
On $\{\sup_{t<T_x}V<K\}$ we have $\rho_n'\uparrow T_x$, so monotone convergence gives
$\int_0^{T_x}\frac{dt}{V_tD_t^2}<\infty$ a.s.\ there. Since $\frac{1}{XX^y}=\frac{1}{D^2V(V+1)}\le\frac{1}{D^2V}$,
\[\log D_{T_x}=\log(y-x)-a\int_0^{T_x}\frac{dt}{X_tX^y_t}>-\infty,\]
i.e.\ $D_{T_x}>0$ and $T_x<T_y$. Letting $K\to\infty$ along $\N$, $T_x<T_y$ a.s.

\end{document}
